% !TEX root = 米粒粘连导致漏数的问题与解决.tex
\appendix
% 附录每节是“一段说明 + 两张图”，就地排版让图紧跟说明，不会被甩到别的小节里
\renewenvironment{cfigure}{\begin{figure}[H]\centering}{\end{figure}}
% 附录的图都用 [H] 就地放，若仍按 \flushbottom 把每页撑满，图与标题之间会被拉出大段空白
\raggedbottom

\section{更多计数结果}
\label{app:renders}

D2 按顺序取两个场景。D1 与 D3 按顺序取到的第一张分别就是正文图 \ref{fig:render_d1} 与图 \ref{fig:render_d3}，
为免重复，D1 改取其余照片中逐粒核对误差最小的一张，D3 改取排在三分之一与三分之二处的两张。
左为本文方法、右为 SAM 3。

\subsection{真实照片 D1}

图 \ref{fig:app_d1} 的场景真值 123 粒，本文方法数出 124 粒，SAM 3 数出 118 粒。
拿标注框逐粒核对，本文方法没有漏数，只在一处没有米的地方多认出一粒。硬币只勾灰边，不计数。

图 \ref{fig:app_d1b} 的场景真值 82 粒，本文方法数出 82 粒，SAM 3 数出 84 粒。
总数一粒不差，逐粒看却有出入。
有 3 处一块区域里其实有两粒米，却只算作一粒；其中最大的一块面积是单粒的 2.2 倍，
凸实度却与单粒几乎相同，没有被判为粘连。
另有 3 块碎片被多算，大小约为单粒的 0.45 倍，恰在碎屑门限附近。
一少一多相抵，总数才恰好等于真值，这也是只看总数看不出来的地方。

\begin{cfigure}
  \includegraphics[width=\textwidth]{appendix_d1_a.png}
  \caption{真实照片 D1 场景一，左为本文方法、右为 SAM 3。}
  \label{fig:app_d1}
\end{cfigure}

\begin{cfigure}
  \includegraphics[width=\textwidth]{appendix_d1_b.png}
  \caption{真实照片 D1 场景二，左为本文方法、右为 SAM 3。}
  \label{fig:app_d1b}
\end{cfigure}

\subsection{合成粘连图 D2}

图 \ref{fig:app_d2} 的场景没有粘连，两种方法都数出 91 粒，和真值一致。
第二个场景粘连率 40\%，真值 93 粒，本文方法 93 粒、SAM 3 94 粒。
图 \ref{fig:app_d2b} 中没有虚线框，粘连块都被真正切开了。

\begin{cfigure}
  \includegraphics[width=\textwidth]{appendix_d2_a.png}
  \caption{合成粘连图 D2 中不粘连的场景，左为本文方法、右为 SAM 3。}
  \label{fig:app_d2}
\end{cfigure}

\begin{cfigure}
  \includegraphics[width=\textwidth]{appendix_d2_b.png}
  \caption{合成粘连图 D2 中粘连率 40\% 的场景，左为本文方法、右为 SAM 3。}
  \label{fig:app_d2b}
\end{cfigure}

\subsection{低分辨率图 D3}

图 \ref{fig:app_d3} 的场景真值 17 粒，本文方法数出 18 粒，SAM 3 数出 17 粒。
多出的一粒在左上方标 $\times 2$ 的块上。那一块面积约为单粒的两倍，又带凹口，被当成两粒粘连，
其实只有一粒米。
图 \ref{fig:app_d3b} 的场景真值 15 粒，两种方法都数出 15 粒，逐粒核对也对得上。
两张图里木框边缘的亮块和阴影都被判为异物，灰色标出，没有计数。
米粒只有 4 像素宽时，形状所能提供的信息已经很少，D3 也是 SAM 3 唯一比本文方法准的一组数据。

\begin{cfigure}
  \includegraphics[width=\textwidth]{appendix_d3_a.png}
  \caption{低分辨率图 D3 场景一，左为本文方法、右为 SAM 3。}
  \label{fig:app_d3}
\end{cfigure}

\begin{cfigure}
  \includegraphics[width=\textwidth]{appendix_d3_b.png}
  \caption{低分辨率图 D3 场景二，左为本文方法、右为 SAM 3。}
  \label{fig:app_d3b}
\end{cfigure}

\section{代码、数据与复现方式}
\label{app:repro}

代码、数据来源说明与复现步骤都放在 \url{https://github.com/stars-spark/asw-spc-rice-counting}。欢迎老师检查！
